\documentclass[aspectratio=169,UTF8]{ctexbeamer}
\usetheme{Madrid}
\usecolortheme{default}
\definecolor{DeepBlue}{HTML}{2563EB}
\definecolor{DeepInk}{HTML}{172033}
\definecolor{DeepGreen}{HTML}{059669}
\setbeamercolor{structure}{fg=DeepBlue}
\setbeamercolor{title}{fg=white,bg=DeepBlue}
\setbeamercolor{frametitle}{fg=white,bg=DeepBlue}
\setbeamercolor{alerted text}{fg=DeepGreen}
\setbeamertemplate{navigation symbols}{}
\setbeamertemplate{footline}[frame number]
\usepackage{booktabs}
\usepackage{listings}
\usepackage{xcolor}
\usepackage{hyperref}
\hypersetup{colorlinks=true,urlcolor=DeepBlue}
\lstset{
  basicstyle=\ttfamily\scriptsize,
  keywordstyle=\color{DeepBlue},
  stringstyle=\color{DeepGreen},
  commentstyle=\color{gray},
  showstringspaces=false,
  breaklines=true,
  frame=single,
  columns=fullflexible
}
\newcommand{\api}[1]{\texttt{#1}}
\newcommand{\keypoint}[1]{\alert{#1}}
\AtBeginSection[]{
  \begin{frame}{本节路线}
    \tableofcontents[currentsection]
  \end{frame}
}

\title{第7课：图像理解}
\subtitle{URL、Base64、OpenCV、文件引用与视觉预算}
\author{DeepSeek API 全栈课程}
\date{文档快照：2026-11-05}
\begin{document}
\frame{\titlepage}
\begin{frame}{学习目标}
\begin{itemize}
  \item 用 URL、Base64、OpenCV 与 file\_id 向模型传图
  \item 选择 detail 级别并估算图片 Token
  \item 掌握格式、大小、尺寸、数量与角色限制
  \item 理解 Chat、Responses、Anthropic 的块结构差异
\end{itemize}
\end{frame}
\section{视觉消息}
\begin{frame}{模型与格式}
仅 \api{deepseek-flash} 处理图像；\api{deepseek-v4-pro} 不支持。支持 JPEG、PNG、GIF、WebP，以文件实际内容而非扩展名判断。\\[8pt]
典型任务：描述、OCR、截图理解、图表分析、视觉问答。
\end{frame}
\begin{frame}{OpenAI Chat 内容块}
user 的 content 从字符串改为数组：
\begin{itemize}
  \item \api{\{"type":"text","text":"..."\}}
  \item \api{\{"type":"image\_url","image\_url":\{...\}\}}
  \item 或 \api{\{"type":"file","file\_id":"..."\}}
\end{itemize}
图片只能出现在 user 消息；system/assistant 携带会返回 400。
\end{frame}
\begin{frame}{三种传图方式}
\begin{tabular}{lll}
\toprule
方式 & 优点 & 适用\\
\midrule
Base64 data URL & 本地文件直接发送 & 小图、一次性\\
公开 HTTP(S) URL & 请求体小 & 稳定公开资源\\
Files API file\_id & 复用、支持更大文件 & 多次请求/大图\\
\bottomrule
\end{tabular}
\end{frame}
\begin{frame}[fragile]{Base64 内容块}
\begin{lstlisting}[language=Python]
b64 = base64.b64encode(path.read_bytes()).decode()
content = [
  {"type": "text", "text": "描述图片"},
  {"type": "image_url", "image_url": {
    "url": f"data:image/jpeg;base64,{b64}",
    "detail": "low"
  }}
]
\end{lstlisting}
内联编码计入 48 MiB 请求体限制。
\end{frame}
\begin{frame}{外部 URL 约束}
URL 最长 8192 字符，图片最大 32 MiB，服务器需在 60 秒内下载完成。URL 必须稳定、可公开访问且返回真正图片。\\[8pt]
私有 URL 不应通过临时放开权限来迁就模型；改用 Base64 或 Files API。
\end{frame}
\section{OpenCV 内存图像}
\begin{frame}{为什么不能直接发送 ndarray？}
\api{cv2.imread()} 返回 NumPy BGR 像素矩阵。\api{image.tobytes()} 只是裸像素，不包含 JPEG/PNG 文件头、尺寸、通道与数据类型，不能作为视觉图片直接放进 JSON。\\[8pt]
正确流程：
\begin{center}
ndarray $\rightarrow$ cv2.imencode $\rightarrow$ 图片格式字节 $\rightarrow$ Base64 或 Files API
\end{center}
\end{frame}
\begin{frame}[fragile]{OpenCV 编码后发送}
\begin{lstlisting}[language=Python]
成功, 缓冲区 = cv2.imencode(
    ".jpg", 图片,
    [cv2.IMWRITE_JPEG_QUALITY, 85],
)
if not 成功:
    raise RuntimeError("图片编码失败")
图片字节 = 缓冲区.tobytes()
Base64文本 = base64.b64encode(图片字节).decode("ascii")
图片地址 = f"data:image/jpeg;base64,{Base64文本}"
\end{lstlisting}
OpenCV 的 BGR 矩阵可直接交给 \api{imencode}，无需先转 RGB。
\end{frame}
\begin{frame}{编码格式与发送路径}
\begin{tabular}{lll}
\toprule
选择 & 建议场景 & 注意事项\\
\midrule
JPEG & 照片、摄像头帧 & 有损；可调质量压缩\\
PNG & OCR、界面、图表 & 无损；通常体积更大\\
Base64 & 单次使用 & 体积约增加三分之一\\
Files API & 多次复用/较大图片 & 上传后需管理生命周期\\
\bottomrule
\end{tabular}
\end{frame}
\section{细节与 Token}
\begin{frame}{detail 级别}
\begin{description}
  \item[low] 缩放到 512$\times$512，更快省 Token
  \item[original] 保留原图
  \item[high] 当前等价于 original
  \item[auto] 当前等价于 original
\end{description}
OCR、小字、密集图表用 original；粗粒度分类/描述先试 low。
\end{frame}
\begin{frame}{图片 Token 规则}
小图保持比例放大到约 544$\times$544；大图缩小到总像素约等于 1300$\times$1300。每张图最多 1024 Token，多图逐张计算。\\[8pt]
实际以响应 usage 为准；超高分辨率不会无限增加模型可见细节。
\end{frame}
\section{完整限制}
\begin{frame}{大小与数量}
\small
\begin{tabular}{ll}
\toprule
限制 & 数值\\
\midrule
请求体 & 48 MiB\\
单图 Base64 / URL & 32 MiB\\
单图 file\_id & 64 MiB\\
单请求图片数 & 600\\
图片总大小 & 无 file\_id 64 MiB；含 file\_id 200 MiB\\
最长边 & 8192 px；15 张及以上时 4096 px\\
\bottomrule
\end{tabular}
\end{frame}
\begin{frame}{跨 API 的块名}
\begin{tabular}{lll}
\toprule
API & 图片块 & 来源字段\\
\midrule
Chat & image\_url / file & url 或 file\_id\\
Responses & input\_image & image\_url 或 file\_id\\
Anthropic & image & source: base64/url/file\\
\bottomrule
\end{tabular}
Anthropic 的 file 来源还需 beta 请求头；第8、10、11课展开。
\end{frame}
\begin{frame}{提示与验证}
\begin{itemize}
  \item 指定关注区域、输出格式与不确定性要求
  \item OCR 结果应回看原图；关键数据做规则校验
  \item 模型可能误读细字、遮挡、比例尺或图例
  \item 人脸、证件、医疗等敏感图片需满足合规要求
\end{itemize}
\end{frame}
\begin{frame}{本课例程}
\begin{enumerate}
  \item \api{01\_vision\_base64.py}：本地图片内联
  \item \api{02\_vision\_url.py}：公开 URL 与 detail
  \item \api{03\_vision\_opencv.py}：OpenCV 内存编码与两种发送路径
\end{enumerate}
第8课继续讲解 file\_id 的查询、分页、过期与删除。
\end{frame}
\begin{frame}{官方资料}
\small
\url{https://api-docs.deepseek.com/zh-cn/guides/vision}\\
\url{https://api-docs.deepseek.com/zh-cn/quick_start/token_usage}\\
\url{https://api-docs.deepseek.com/zh-cn/api/create-chat-completion}
\end{frame}
\end{document}
